% Generated draft: requires mathematical and source review.
\documentclass[11pt]{article}
\usepackage[margin=25mm]{geometry}
\usepackage{fontspec}
\setmainfont{Latin Modern Roman}
\usepackage{amsmath,amssymb,mathtools,unicode-math}
\setmathfont{Latin Modern Math}
\usepackage{enumitem,needspace,xurl,hyperref}
\setlength{\parindent}{0pt}
\setlength{\parskip}{5pt}
\setlength{\emergencystretch}{3em}
\newcommand{\N}{\mathbb{N}}
\newcommand{\Z}{\mathbb{Z}}
\newcommand{\Q}{\mathbb{Q}}
\newcommand{\R}{\mathbb{R}}
\newcommand{\C}{\mathbb{C}}
\newcommand{\F}{\mathbb{F}}
\newcommand{\A}{\mathbb{A}}
\newcommand{\PP}{\mathbb{P}}
\newcommand{\E}{\mathbb{E}}
\newcommand{\eps}{\varepsilon}
\newcommand{\dd}{\,\mathrm{d}}
\begin{document}
\section*{UnsolvedMath ID: 30005510}
\textbf{Source problem code:} catalogue entry
\subsection{Definitions and mathematical statement}
Let \(G\) be a finite group, and let \(BG\) denote its classifying space. The supplied record uses \(C^*BG\) for a cochain algebra on \(BG\), \(D_{\mathrm{sg}}(C^*BG)\) for its singularity category, and \(\Theta_G\) for the group-theoretic nucleus. It does not specify the coefficients, the construction of this singularity category, the support theory, or the definition of the nucleus.

With these constructions understood in the intended conventions of the cited report, the recorded conjecture is
\[
\operatorname{supp}\!\left(D_{\mathrm{sg}}(C^*BG)\right)=\Theta_G
\qquad\text{for every finite group }G.
\]
This is a statement of the supplied target, subject to unresolved definitions. A fully specified formulation requires the coefficient assumptions, the module categories and quotient defining \(D_{\mathrm{sg}}\), the ambient space and meaning of support, and the construction of \(\Theta_G\). Any dependence of these objects on the coefficient characteristic must also be made explicit. No coefficient domain or support convention is imposed here.

The supplied original statement reads, ``We conjecture that this is the case in general.'' Its antecedent is absent from the supplied passage, so that sentence alone does not determine the displayed equality or its hypotheses.
\subsection{Short English statement}
The recorded conjecture says that, for every finite group \(G\), the support of the singularity category of cochains on \(BG\) equals the group-theoretic nucleus \(\Theta_G\). The supplied record omits the coefficient assumptions and definitions needed to make this equality precise. It labels the problem open; that label is supplied curation and has not been independently verified.
\subsection{Sources}
\begin{itemize}
\item Supplied local UnsolvedMath registry record, permanent ID \(30005510\), title: Singularity-Category Support and the Group-Theoretic Nucleus. Exact source problem code: \texttt{catalogue entry}. The record assigns the problem to 2023 and supplies no named proposer.
\item Report provenance supplied by the registry: \emph{Representations of Finite Groups} (2023), Oberwolfach Reports, DOI 10.4171/owr/2023/19, pp.\ 1031--1084. \url{https://doi.org/10.4171/owr/2023/19}. The background describes Oberwolfach Reports as records of mathematical workshops and their open-problem discussions.
\item Supplied literature source: Benson--Carlson, \emph{Modules with finitely generated cohomology}. \url{https://doi.org/10.1016/j.jalgebra.2024.07.040}.
\item Supplied literature source: Benson--Greenlees, \emph{modules and singularities of C-star-BG}. \url{https://arxiv.org/abs/2305.08580}.
\item Supplied literature source: \emph{Oberwolfach Report 19/2023}. \url{https://doi.org/10.4171/owr/2023/19}.
\item Literature assessment dated exactly \texttt{2026-08-22}: the supplied curation labels the problem open and states that the equality is proved in special cases, including additional examples from 2024, while the general equality remains unproved. This is supplied dated curation, not an independently verified literature assessment. The background requests independent verification of time-sensitive status and cited claims.
\item The registry labels the statement \texttt{corrected\_verified} and reports that it was reconstructed and checked against the official Oberwolfach report. That verification claim is supplied metadata; it has not been independently checked for this draft.
\end{itemize}
\end{document}
